\chapter{What These Chapters Owe}
\label{sec:hyper-gaps}
% ===================================================================

Before the act turns to reference, an accounting.
The chapters just finished make the loudest claims in the book and rest
on the thinnest construction in it, and
Remark~\ref{rmk:why-here} promised that this arrangement would be stated
rather than disguised.
This is the statement.

\section{Open gaps}

\begin{enumerate}[label=\textbf{Gap \arabic*.}, leftmargin=*, itemsep=8pt]

  \item \textbf{Colimit at limit ordinals.}
        Definition~\ref{def:oracle-ext} defines $\fiber{\lambda}$ as the
        colimit in $\mathbf{GSLT}$ for limit ordinals $\lambda$.
        The existence and properties of this colimit require
        verification.
        In particular: does the colimit GSLT have a well-defined
        bisimulation quotient, and is it the union of the bisimulation
        quotients of the stages below?

  \item \textbf{Subject reduction for schedulers.}
        This gap originally asked for a construction of the oracle as a
        persistent receive in the rho calculus.
        Chapters~\ref{ch:apr} and~\ref{ch:agy} argue the gap was the
        wrong shape: the oracle need not be constructed at all, because
        the resolution of a contended channel is already a choice
        function and the scheduler that supplies it already spends a
        graded choice resource.
        What is owed instead is the metatheorem.
        Does the term assignment of Chapter~\ref{ch:chs} satisfy subject
        reduction --- can a scheduler, merely by running, come to realize
        a stronger choice principle than it was typed for?
        If it can, the graded factorization is not a reduction invariant,
        an island can crack while it runs, and that cracking is what one
        agent perturbing another amounts to.
        This is Question~\ref{ndo:q:crux} and Question~\ref{ndo:q:sr},
        which are the same question, and it is the load-bearing gap of
        these chapters.

  \item \textbf{Strict refinement over the lattice order
        (Proposition~\ref{prop:strict-bisim}).}
        The proof sketch uses a diagonalization argument, and it was
        restated in this pass for the lattice order $v < w$ rather than
        for an ordinal successor.
        The full proof requires specifying the encoding of the
        scheduling problem characteristic of $w$ as terms and verifying
        that the diagonalization goes through in the GSLT setting.
        The restatement makes the gap slightly wider than it was: for a
        successor there was an obvious candidate problem, and for a
        general pair $v < w$ one must be produced.

  \item \textbf{New currents (Proposition~\ref{prop:new-noether}).}
        The claim that new conserved quantities appear higher in the
        lattice requires identifying concrete symmetries of the cost map
        at $w$ and verifying that they give rise to conserved quantities
        via the Noether argument of Part~\ref{part:physics}.

  \item \textbf{The standing charge of a position.}
        Section~\ref{sec:relationship} asserts that occupying a lattice
        point is a recurring cost and not a one-time acquisition, and the
        unfavorable cycle turns on it.
        Chapter~\ref{ch:agy} prices individual resolutions; it does not
        price the maintenance of the capacity to resolve.
        Nothing in the book supplies that quantity, and
        Remark~\ref{rmk:descent} needs it.

  \item \textbf{Effective resolution.}
        Definition~\ref{def:consciousness} requires that the agent can
        \emph{form and evaluate} the queries a position resolves.
        A precise definition of effective access, in terms of the agent's
        account and the expressiveness of its hypothesis language, is
        needed.

  \item \textbf{Approximate resolution
        (Definition~\ref{def:approx}).}
        Remark~\ref{rmk:approx-status} restates this on the stochastic
        instance, correcting a squared amplitude that should not have
        survived the withdrawal in Chapter~\ref{sec:quantum}.
        What the restatement leaves open is whether approximate
        resolution places an agent at a genuine intermediate point of
        $\mathfrak{W}$, or merely grades its reliability at a point it
        already occupies.
        These are different claims and the earlier presentation ran them
        together.

  \item \textbf{Internalization is a conjecture
        (Conjecture~\ref{suf:conj:selected}).}
        Chapter~\ref{ch:suffer} states the selection argument for
        internalization in the form it would need to be checked in ---
        a threshold in the variance of inflow --- and does not check it.
        The machinery to check it exists: Chapter~\ref{ch:weight}
        supplies propensities and a well-posed simulator.
        The experiment has not been run, and until it is, the price of a
        self-model is located rather than known.

  \item \textbf{The hard problem.}
        These chapters offer structural analogues: mathematical objects
        with the right functional properties.
        They do not address why there is something it is like to be an
        agent reading its own drift at a point in the lattice.
        The explanatory gap between functional description and phenomenal
        experience remains open and is not claimed to be resolved here.
        Chapter~\ref{ch:suffer} chooses a heavily freighted word for a
        structure, deliberately; this gap is the price of that choice and
        we do not pretend the word pays it.

\end{enumerate}

\section{Relation to existing work}
\label{sec:hyper-discussion}

The book's comparisons with neighboring programs --- integrated
information theory, the free energy principle, global workspace,
constructor theory, the ruliad, AIXI, and assembly theory --- are gathered
in Chapter~\ref{ch:related} at the end of the Pledge, where a reader can
find them before committing to six hundred pages rather than after.
What belongs here is only what is specific to these chapters.

Turing's oracle hierarchy~\cite{turing1939} is the direct inspiration for
the tower of Chapter~\ref{ch:tower}, and the departure is the move off
the chain: first to a family indexed over a category of theories, and
then, once Remark~\ref{rmk:not-a-fibration} has conceded that this family
is not the fibration it was called, to a coordinate in the Weihrauch
lattice.

The account of sentience is where the largest change from earlier drafts
lies, and it is worth saying what it was.
Sentience was a distinguished component of the vectorial account driven
by a reinforcement signal on predictive accuracy, which made it formally
a reinforcement-learning construct with an internal
reward~\cite{sutton2018reinforcement}.
Chapter~\ref{ch:suffer} withdraws that and replaces it with a formula the
learner holds about its own reserves.
The gain is that the reward function stops being a free parameter: there
is nothing left to choose, because the signal is the learner's own assay
of its distance from a goal it carries.
The loss is that the construction now depends on the instrument
discipline of Chapter~\ref{ch:obsext}, which is itself conditional on two
unproved hypotheses.
We think that is the right trade --- a derived quantity resting on stated
hypotheses is better than a posited one resting on nothing --- but the
reader should see it as a trade.

One correction, since earlier drafts made the claim and it should not
stand.
The treatment of consciousness was described as ordinal-valued.
That description was superseded by Chapters~\ref{ch:apr}
and~\ref{ch:agy}, which locate the vertical coordinate in the Weihrauch
lattice --- emphatically a lattice and not a chain.
Consciousness here is a position an agent occupies, not a threshold it
crosses, and the ordinal tower survives only as the chain-shaped shadow
used to introduce the intuition.

\section{Where the argument stands}

Taken together with what precedes them, these chapters complete the
framework's third leg.

\begin{itemize}[itemsep=4pt]
  \item Part~\ref{part:physics} establishes the physics: causal
        structure, dynamics, and a semiring-parametric path weight from
        the bare structure of a GSLT.
        The complex instance of that weight was proposed and then
        withdrawn in Chapter~\ref{sec:quantum}; what the framework
        carries is the Boolean, tropical and stochastic instances, and
        the last of those is what everything downstream actually uses.

  \item Part~\ref{part:sciencebot} establishes the epistemology: why
        situated agents in massive populations perceive a hierarchical
        false ontology rather than the true bisimulation quotient.

  \item These chapters establish the phenomenology, in the deflationary
        sense the book can support: why a learner that watches its own
        reserves has states that behave like affect, and why a learner's
        position fixes not only what it can do but what there is.
\end{itemize}

The unifying theme is that the category of GSLTs is not merely a setting
for studying computation.
It is a setting rich enough to ground physics, epistemology, and --- if
the arguments here have merit --- phenomenology.

What remains is the question the Prestige exists to ask.
All of the above was got by restricting scope to computation, and the
restriction bought an ontology and a grounded language free.
Chapter~\ref{ch:notation} asks what would be required to ground a symbol
system without that restriction, and Chapter~\ref{ch:sim} asks what the
answer implies about the world we are in.

% ===================================================================
